\subsection{位与赋值环}

命题 1. 设 K 是域，A 是 K 的赋值环，\(\kappa(A)\) 是 A 的剩余域。把从 A 到 \(\kappa(A)\) 上的典范映射扩张为映射 \(h_{A}: \tilde{K} \to (\kappa(A))^{\sim}\)，依据等式 \(h_{A}(x) = \infty\)（当 \(x \notin A\)）。如此定义的映射 \(h_{A}\) 是 K 的一个位，其值域为 \(\kappa(A)\)。

显然 \(h_{\mathbf{A}}(1) = 1\)。

我们证明 \(h_{\mathbf{A}}\) 是关于加法的态射。设 \(x, y\) 是 \(\tilde{\mathbf{K}}\) 的两个元素，使 \(h_{\mathbf{A}}(x) + h_{\mathbf{A}}(y)\) 有定义。这时两个元素 \(x, y\) 之一属于 \(\mathbf{A}\)，从而 \(x + y\) 有定义。若 \(x \in \mathbf{A}\) 且 \(y \in \mathbf{A}\)，显然

\[
h _ {\mathrm{A}} (x) + h _ {\mathrm{A}} (y) = h _ {\mathrm{A}} (x + y)
\]

成立。若 \(x \in \mathbf{A}\) 而 \(y \notin \mathbf{A}\)，则 \(x + y \notin \mathbf{A}\)，上式的两端都等于 \(\infty\)。

最后证明 \(h_{\mathbf{A}}\) 是关于乘法的态射。设 \(x \in \tilde{\mathbf{K}}, y \in \tilde{\mathbf{K}}\) 使 \(h_{\mathbf{A}}(x)h_{\mathbf{A}}(y)\) 有定义。若 \(x \in \mathbf{A}\) 且 \(y \in \mathbf{A}\)，显然 \(xy\) 有定义且 \(h_{\mathbf{A}}(x)h_{\mathbf{A}}(y) = h_{\mathbf{A}}(xy)\)。现在设元素 \(x, y\) 之一（例如 \(y\)）不属于 \(\mathbf{A}\)；由于 \(h_{\mathbf{A}}(y) = \infty\)，有 \(h_{\mathbf{A}}(x) \neq 0\)，即 \(x \notin \mathfrak{m}(\mathbf{A})\)，从而 \(x^{-1} \in \mathbf{A}\)；由此得出 \(xy\) 有定义且 \(xy \notin \mathbf{A}\)，于是

\[
h _ {\mathrm{A}} (x y) = \infty = h _ {\mathrm{A}} (x) h _ {\mathrm{A}} (y).
\]

这就证明了命题 1。

若 \(j\) 是从 \(\kappa(A)\) 到某个域 \(L, j \circ h_{A}: \tilde{K} \to \tilde{L}\) 的一个子域上的同构，则它是 \(K\) 的以 \(L\) 为值域的位。上述过程实际上给出了 \(K\) 上的所有位。确切地说：

命题 2. 设 K、L 是两个域，f 是 K 的以 L 为值域的位。则存在 K 的赋值环 A 与从 \(\kappa(\mathbf{A})\) 到 L 的一个子域上的同构 j，使 \(f = j \circ h_{A}\)；这些条件唯一确定 A 与 j。环 A 是由 \(x \in K\) 中使 \(f(x) \neq \infty\) 的元素组成之集，\(\mathfrak{m}(\mathbf{A})\) 是由 \(x \in K\) 中使 \(f(x) = 0\) 的元素组成之集。

若 \(f = j \circ h_{\mathbf{A}}\)，则条件 \(f(x) \neq \infty\)（相应地，\(f(x) = 0\)）等价于条件 \(h_{\mathbf{A}}(x) \neq \infty\)（相应地，\(h_{\mathbf{A}}(x) = 0\)），从而等价于条件 \(x \in \mathbf{A}\)（相应地，\(x \in \mathfrak{m}(\mathbf{A})\)）。于是 \(\mathbf{A}\) 被唯一确定，且由于 \(h_{\mathbf{A}}\) 是满射，\(j\) 也是唯一的。

现在设 \(f\) 是 \(\mathbf{K}\) 的以 \(\mathbf{L}\) 为值域的任一位；用 \(\mathbf{A}\) 表示由 \(x \in \mathbf{K}\) 中使 \(f(x) \neq \infty\) 的元素组成之集。若 \(x \in \mathbf{A}\) 且 \(y \in \mathbf{A}\)，则合成 \(f(x) - f(y)\) 与 \(f(x)f(y)\) 有定义且 \(\neq \infty\)，这表明 \(x - y \in \mathbf{A}\) 且 \(xy \in \mathbf{A}\)；于是 \(\mathbf{A}\) 是 \(\mathbf{K}\) 的子环。若 \(x \notin \mathbf{A}\)，则 \(f(x) = \infty\)，从而 \(f(x^{-1}) = 0\)，于是 \(x^{-1}\) 属于核 \(m\)，其中 \(f'\) 是把 \(f\) 限制到 \(\mathbf{A}\) 所得的同态。反之，若 \(y \in m\)，则 \(y^{-1} \notin \mathbf{A}\)。这表明 \(\mathbf{A}\) 是 \(\mathbf{K}\) 的赋值环，且 \(m\) 是它的极大理想。设 \(j\) 是从 \(\kappa(\mathbf{A})\) 到 \(L\) 的单同态，它是由 \(f'\) 通过取商导出的。则 \(f(x) = j(h_{\mathrm{A}}(x))\)（对所有 \(x \in A\)），且这个等式对 \(x \notin A\) 仍成立，此时两端都等于 \(\infty\)。

分解 \(f = j \circ h_{A}\) 称为位 f 的典范分解。A 称为 f 的环，\(\mathfrak{m}(A)\) 称为 f 的理想，\(\kappa(A)\) 称为 f 的剩余域。为使两个位 \(f: \tilde{K} \to \tilde{L}\) 与 \(f': \tilde{K} \to \tilde{L}'\) 有相同的环，必须且只需存在从 f 的值域到 \(f'\) 的值域上的同构 s，使 \(f' = s \circ f\)；这时称 f 与 \(f'\) 等价。可以看出，关于赋值环的每个结果都可以翻译成关于位的结果，反之亦然；下面各节我们要做的正是这件事。

\textbf{位的例子}\par

\begin{enumerate}
\def\labelenumi{(\arabic{enumi})}
\item
  设 \(\mathbf{K}\) 是域。\(\mathbf{K}\) 上的恒等映射是平凡位，其环为 \(\mathbf{K}\)，理想为 (0)。
\item
  设 \(k\) 是域。对所有 \(u \in k((\mathrm{T}))^{\sim}\)，记 \(f(u) = \infty\)（当 \(u \notin k[[\mathrm{T}]]\) 时），并定义 \(f(u)\) 为 \(u\) 的常数项（当 \(u \in k[[\mathrm{T}]]\) 时）。则 \(f\) 是 \(k((\mathrm{T}))\) 的一个位，其剩余域为 \(k\)，环为 \(k[[\mathrm{T}]]\)。因为 \(k[[\mathrm{T}]]\) 是 \(k((\mathrm{T}))\) 的赋值环（§ 1 第 4 节例 3），且 \(f\) 在 \(k[[\mathrm{T}]]\) 上的限制等同于从 \(k[[\mathrm{T}]]\) 到它的剩余域上的典范同态。
\item
  设 \(k\) 是域，\(a\) 是 \(k\) 的元素，\(A\) 是由这样的 \(u \in k(X)\) 组成之集：\(a\) 可代入 \(u\)（《代数》第四章 § 3 第 2 节）。若用 \(p\) 表示素理想 \((X - a)\)（\(k[X]\) 的），则 \(A = k[X]_p\)，从而 \(A\) 是 \(k(X)\) 的赋值环（§ 1 第 4 节命题 2）。对所有 \(u \in k(X)^{\sim}\)，记 \(f(u) = \infty\)（当 \(u \notin A\) 时），记 \(f(u) = u(a)\)（当 \(u \in A\) 时）。则 \(f\) 是 \(k(X)\) 的一个位，其剩余域为 \(k\)，环为 \(A\)；因为 \(f\) 在 \(A\) 上的限制是 \(A\) 到 \(k\) 上的同态（《代数》第四章 § 3 命题 2），其核为 \(pA = m(A)\)。称元素 \(f(u) \in \tilde{k}\) 是由令 \(X = a\)（在 \(u\) 中）而得到的。
\end{enumerate}

* (4) 设 S 是一维连通复解析簇，K 是 S 上的亚纯函数域。对所有 \(z_0 \in \mathbf{S}\)，映射 \(f \mapsto f(z_0)\)（从 K 到 \(\tilde{\mathbf{C}}\)）是 K 的一个位，其环是 \(f \in \mathbf{K}\) 中在 \(z_0\) 处全纯的那些元素所成之集，其理想是 \(f \in \mathbf{K}\) 中在 \(z_0\) 处为零的那些元素所成之集。正是这个例子以及其他类似的例子，成为术语“位”的来源。*

